\documentclass{beamer}

\usepackage{amsmath,amsfonts,amssymb,bm}
\usepackage{braket}
\usepackage{physics}

\usetheme{Madrid}

\title{How the HHL Algorithm is Used in Quantum Machine Learning}
\subtitle{Linear Systems, Kernels, and Spectral Filtering}
\author{Morten Hjorth-Jensen}
\date{Spring 2026}

\begin{document}

\frame{\titlepage}

%================================================
\section{Motivation}
%================================================

\begin{frame}{The Core Idea}
The HHL algorithm solves a linear system
\[
A x = b
\]
by preparing a quantum state proportional to
\[
\ket{x} \propto A^{-1}\ket{b}.
\]

In quantum machine learning, this is important because many learning tasks reduce to
\begin{itemize}
\item linear systems,
\item matrix inversion,
\item least-squares problems,
\item kernel methods.
\end{itemize}
\end{frame}

%------------------------------------------------

\begin{frame}{Why HHL Appears in Machine Learning}
Many machine learning algorithms can be reformulated as
\[
A x = b
\]
for a suitable matrix \(A\) and vector \(b\).

Examples include:
\begin{itemize}
\item linear regression,
\item ridge regression,
\item kernel methods,
\item Gaussian processes,
\item linearized neural networks.
\end{itemize}

Thus HHL serves as a quantum subroutine for solving the underlying learning problem.
\end{frame}

%================================================
\section{Quantum Linear Regression}
%================================================

\begin{frame}{Ordinary Least Squares}
In linear regression, one minimizes
\[
\min_w \norm{Xw - y}^2.
\]

The formal solution is
\[
w = (X^T X)^{-1} X^T y.
\]

This can be written as a linear system:
\[
A w = b,
\qquad
A = X^T X,
\qquad
b = X^T y.
\]
\end{frame}

%------------------------------------------------

\begin{frame}{How HHL Enters}
Instead of returning the vector \(w\) explicitly, HHL prepares a quantum state proportional to
\[
\ket{w} \propto (X^T X)^{-1} X^T \ket{y}.
\]

This means:
\begin{itemize}
\item HHL does not directly output all regression coefficients,
\item but it gives quantum access to the solution state,
\item which can be used to estimate global properties or predictions.
\end{itemize}
\end{frame}

%================================================
\section{Kernel Methods and QSVMs}
%================================================

\begin{frame}{Kernel Ridge Regression and Related Methods}
Kernel methods often require solving
\[
(K + \lambda I)\alpha = y,
\]
where:
\begin{itemize}
\item \(K\) is the kernel matrix,
\item \(\lambda\) is a regularization parameter,
\item \(\alpha\) determines the classifier or regressor.
\end{itemize}

This is again a linear system, so HHL can in principle be applied.
\end{frame}

%------------------------------------------------

\begin{frame}{Connection to Quantum Support Vector Machines}
In a quantum kernel method, one defines
\[
K(x,x') = |\braket{\phi(x)}{\phi(x')}|^2,
\]
where \(\ket{\phi(x)}\) is a quantum feature state.

A more fully quantum workflow would:
\begin{enumerate}
\item compute the quantum kernel entries on a quantum computer,
\item use HHL to solve the linear system involving \(K\),
\item use the resulting coefficients for classification or regression.
\end{enumerate}
\end{frame}

%================================================
\section{Gaussian Processes}
%================================================

\begin{frame}{Gaussian Process Regression}
Gaussian process regression requires solving systems of the form
\[
(K + \sigma^2 I)^{-1} y,
\]
where:
\begin{itemize}
\item \(K\) is the covariance or kernel matrix,
\item \(\sigma^2\) is the noise variance.
\end{itemize}

The predictive mean is typically
\[
\mu(x_*) = k_*^T (K + \sigma^2 I)^{-1} y.
\]

This again fits naturally into the HHL framework.
\end{frame}

%================================================
\section{Quantum Neural Networks in the Linearized Regime}
%================================================

\begin{frame}{Neural Tangent Kernel Perspective}
In the linearized regime of neural network training, one approximates
\[
f(x) \approx f(x_0) + \nabla_\theta f \cdot (\theta - \theta_0).
\]

Training then reduces to solving a linear system involving the neural tangent kernel (NTK).

Therefore, HHL could in principle be used as a quantum linear solver in this regime.
\end{frame}

%================================================
\section{Physics Interpretation}
%================================================

\begin{frame}{HHL as a Spectral Inverse}
From a physics perspective, HHL implements
\[
A^{-1}.
\]

If
\[
A = \sum_j \lambda_j \ket{u_j}\bra{u_j},
\]
then
\[
A^{-1} = \sum_j \frac{1}{\lambda_j}\ket{u_j}\bra{u_j}.
\]

So HHL acts as a spectral filter:
\[
\lambda_j \mapsto \frac{1}{\lambda_j}.
\]
\end{frame}

%------------------------------------------------

\begin{frame}{Connection to Green's Functions}
In many-body physics, inverse operators appear as resolvents or Green's functions:
\[
G(z) = (zI - H)^{-1}.
\]

This means that HHL can be understood as applying a Green's-function-like operator to data.

From this viewpoint:
\begin{itemize}
\item regression corresponds to a response problem,
\item kernels resemble propagators,
\item learning becomes spectral filtering.
\end{itemize}
\end{frame}

%================================================
\section{Advantages and Caveats}
%================================================

\begin{frame}{When HHL Can Help}
In principle, HHL can offer strong speedups if:
\begin{itemize}
\item the matrix is sparse or efficiently block-encoded,
\item the condition number is not too large,
\item the input state can be prepared efficiently,
\item one only needs global properties of the solution.
\end{itemize}
\end{frame}

%------------------------------------------------

\begin{frame}{Important Caveats}
There are important limitations:
\begin{itemize}
\item HHL outputs a quantum state, not the full classical vector,
\item reading out all coefficients may destroy the speedup,
\item data loading can be costly,
\item conditioning and noise may limit practical use.
\end{itemize}

Thus HHL is most useful when one wants
\begin{itemize}
\item expectation values,
\item overlaps,
\item predictions,
\item or other global observables.
\end{itemize}
\end{frame}

%================================================
\section{Summary}
%================================================

\begin{frame}{Summary}
HHL enters quantum machine learning as a quantum linear solver for problems such as:
\begin{itemize}
\item linear regression,
\item kernel methods,
\item Gaussian processes,
\item linearized neural networks.
\end{itemize}

Conceptually, one may summarize the structure as
\[
\text{machine learning}
\longleftrightarrow
\text{linear systems}
\longleftrightarrow
\text{resolvents}
\longleftrightarrow
\text{HHL}.
\]

Thus HHL is best viewed as a foundational inverse-operator subroutine in quantum machine learning.
\end{frame}

\end{document}
